\documentclass[11pt]{article}
\usepackage[margin=1in]{geometry}
\usepackage{graphicx,amsmath,amssymb,booktabs,caption,authblk,lineno,microtype}
\usepackage[T1]{fontenc}
\usepackage[hidelinks]{hyperref}
\captionsetup{font=small,labelfont=bf}
\renewcommand{\arraystretch}{1.12}
\linenumbers

\title{\bfseries Malignant-cell callers assign tumour labels to tumour-free tissue:
a null-control audit and a calibrated decision procedure}

\author[1]{[AUTHOR NAMES TO BE SUPPLIED BEFORE SUBMISSION]}
\affil[1]{Affiliation}
\date{}

\begin{document}
\maketitle

\begin{abstract}
\noindent
Single-cell studies of cancer routinely partition cells into malignant and
non-malignant compartments using copy-number inference or supervised
classifiers, and every downstream claim --- tumour heterogeneity, clonal
architecture, malignant-cell programmes --- inherits whatever errors that
partition makes. These methods are almost never evaluated on tissue that
contains no tumour, so their false-positive rate is unknown. We assembled
14 tumour-free single-cell RNA-seq datasets (36{,}528 cells, 76 donors, 10
healthy organs and 4 tissues from inflammatory-disease cohorts) and ran three
widely used malignant-cell callers under five published decision rules. Every
positive call in this panel is, by construction, false. Published rules
produced false malignant calls in every tissue examined, at rates spanning
0.5--19.3\% for the two-criterion rule of Tirosh and colleagues~\cite{tirosh2016} and
4.5--78.0\% for a forced two-group split. CopyKAT called 3.0--63.0\% of cells
aneuploid, exceeding 20\% in 9 of 14 datasets; the supervised classifier
ikarus called 0--100\%, exceeding 50\% in 4 of 14. Pooled over the panel,
CopyKAT and ikarus each assigned roughly 31\% of tumour-free cells to the
malignant class. We then built in-silico positive controls by imposing arm-level
dosage changes of known magnitude on real cells, and found that per-cell
discrimination is near chance for single events (AUC 0.53) and reaches only
0.85 at 18 arm events. At matched false-positive rate, simply ranking cells
by the continuous copy-number score detects roughly twice as many aneuploid
cells as the published rule operating at that rate: the rules do not add
information to the score, they choose a poor operating point on it. The false
calls are not random. They share a reproducible pseudo-karyotype that
correlates across independent tissues (median Spearman $\rho = 0.54$) and is
predicted by local gene density ($\rho = -0.53$, $P = 2\times10^{-21}$):
gene-dense regions appear deleted and gene deserts appear gained. Finally, we
show that a cell-type-conditional empirical null, fitted on tumour-free
reference cells and applied as a calibrated decision rule, holds the realized
false-positive rate at or below its nominal target in 13 of 14 datasets
(median 0.03\%, maximum 0.43\% at $\alpha = 5\%$), and that the single failure
--- a dataset from a different protocol --- is repaired by drawing the null
from held-out donors of the same dataset (27.1\% $\rightarrow$ 0.03\%). We
provide the calibration as a small, caller-agnostic layer that consumes any
per-cell copy-number score.
\end{abstract}

\section{Introduction}

The separation of malignant from non-malignant cells is the first
irreversible decision in most single-cell cancer analyses. It defines the
denominator for tumour-cell abundance, the population in which intratumoural
heterogeneity is measured, the cells from which malignant expression
programmes are derived, and the input to clonal reconstruction. Two families
of method dominate. The first infers large-scale copy-number variation from
average expression in sliding genomic windows, following the approach
introduced for oligodendroglioma, melanoma and
glioblastoma~\cite{tirosh2016b,tirosh2016,patel2014} and implemented in
inferCNV~\cite{infercnv}, CopyKAT~\cite{copykat} and related
tools~\cite{conics,casper,honeybadger,numbat,scevan}; cells are then declared malignant by thresholding
a per-cell summary of the inferred profile, or by clustering cells on that
profile and designating one or more clusters as tumour. The second trains a
supervised classifier on labelled tumour and normal cells and applies it to
new data; ikarus~\cite{ikarus} is the most widely used example.

Both families are evaluated, when they are evaluated at all, on tumour
samples --- where the ground truth is itself often derived from the same or a
closely related signal, and where a cell called malignant cannot easily be
proven normal. The complementary experiment is rarely done: run the caller on
tissue that contains no malignant cells and count the calls. Every such call
is a false positive with certainty, requiring no ground-truth annotation
beyond the provenance of the sample. This is the experiment we report here.

The question matters beyond method benchmarking. A false-positive rate of a
few per cent would be a minor nuisance; a rate of tens of per cent would mean
that a substantial fraction of the cells contributing to published malignant
programmes are normal cells, and that apparent tumour heterogeneity partly
reflects normal cell-type diversity misassigned to the tumour compartment.
The direction of the resulting bias is not random if the false calls are
themselves structured --- if, for instance, the cells most likely to be
miscalled are those with particular transcriptional or technical properties.

We address three questions. First, what is the false-positive rate of current
malignant-cell calling in tumour-free tissue, and how does it depend on the
decision rule rather than on the underlying score? Second, what drives the
false calls --- is there a systematic artefact, and can it be attributed to a
property of the genome or of the assay? Third, can the decision step be
replaced by a calibrated procedure with a controlled and verifiable
false-positive rate, without retraining or replacing the underlying caller?

\section{Results}

\subsection{A tumour-free panel on which every positive call is false}

We assembled 14 single-cell RNA-seq datasets in which no malignant cells are
expected (Supplementary Table S1). Ten are healthy adult organs from a
multi-organ reference atlas~\cite{tabula} --- bone marrow, kidney, large intestine, liver,
lung, pancreas, skin, small intestine, spleen and stomach --- profiled on a
common 10x 3$'$ v3 protocol. Four are drawn from inflammatory-disease
cohorts: epithelial and immune compartments of Crohn's disease
colon~\cite{kong2023}, psoriatic skin~\cite{gao2021pso}, and colon from an
ulcerative colitis cohort~\cite{smillie2019}. The
inflammatory-cohort tissues were included deliberately: inflammation induces
strong, coordinated transcriptional programmes that are a plausible source of
copy-number-like artefacts, and inflamed tissue is a common benign comparator
in cancer studies. After subsampling, the panel comprises 36{,}528 cells from
76 donors.

We ran three callers across the full panel. The first is infercnvpy 0.6.1, a
Python implementation of the inferCNV sliding-window algorithm (100-gene
windows, 10-gene step), run in two reference modes: reference-free (each cell
compared to the dataset mean) and immune-reference (T, B, NK, myeloid and
related cells designated as diploid). The second is CopyKAT 1.2.5, run with
the same immune reference. The third is ikarus 0.0.3, applied with its
published pan-cancer gene signature and classifier. On the
copy-number profiles we applied five decision rules taken from published
practice: a forced two-group split, designation of the highest-scoring
cluster as tumour, a per-cell quantile cutoff, a cutoff at a fixed multiple
of the reference standard deviation, and the two-criterion rule of Tirosh and
colleagues requiring both a high copy-number score and high correlation to
the mean profile of the top-scoring cells.

Every caller and every rule produced false malignant calls, in every tissue
(Figure~\ref{fig:fp}a). With an immune reference, the two-criterion rule
called 0.5--19.3\% of cells malignant depending on the dataset; the forced
two-group split called 4.5--78.0\%. CopyKAT called 3.0--63.0\% of cells
aneuploid, exceeding 20\% in 9 of 14 datasets and reaching 63.0\% in both
healthy large and small intestine, 50.3\% in healthy skin and 49.8\% in
healthy pancreas; pooled over the panel it called 30.7\% of tumour-free cells
aneuploid. A further 0--80.6\% of cells per dataset were returned as
\emph{not.defined} and are excluded from the denominators above; restricting
to cells CopyKAT classified, the pooled rate is 35.3\%. ikarus called
0--100\% of cells tumour, exceeding 50\% in 4 of 14 datasets, with 100\% of
Crohn's colonic epithelium and 91.8\% of healthy liver assigned to the tumour
class, and 30.6\% pooled. The three callers do not agree on which tissues are
problematic: healthy liver is 91.8\% tumour by ikarus and 10.3\% aneuploid by
CopyKAT, while healthy small intestine is 63.0\% aneuploid by CopyKAT and
0.0\% tumour by ikarus.

Two features of this result deserve emphasis. First, the rate is a property
of the decision rule more than of the tissue: applied to the same copy-number
profiles, the rules differ by more than an order of magnitude in pooled
false-call rate (Figure~\ref{fig:fp}b). Second, the score being thresholded
is unimodal (Figure~\ref{fig:fp}c). There is no malignant mode for a
threshold to separate; any rule that must return a positive fraction will cut
into the body of a normal distribution, and the fraction it returns is set by
where the cut falls, not by the presence of tumour.

\subsection{Positive controls: what the callers can actually detect}

A high false-positive rate is only interpretable alongside the sensitivity
that buys it. Because tumour-free tissue provides no positives, we
constructed them: from each of three healthy tissues (lung, liver, spleen) we
drew cells at random and imposed arm-level dosage changes of known burden on
their expression, sampling whole chromosome arms and scaling the counts of
all genes on the affected arms up or down before re-normalisation. This
yields cells that are real in every respect except for a known, controlled
karyotype, alongside untouched diploid companion cells from the same sample
and protocol. We generated 8{,}994 such cells spanning burdens of 1 to 18
arm-level events.

Per-cell discrimination is weak at low burden and moderate at high burden
(Figure~\ref{fig:spike}a). The area under the ROC curve for separating
aneuploid from diploid cells using the continuous copy-number score is 0.53
for a single arm event --- indistinguishable from chance --- rising to 0.85
at 18 events. This sets a ceiling on what any decision rule built on this
score can achieve, and it is a modest ceiling: cells carrying fewer than
about five arm-level changes are not reliably separable from normal cells by
single-cell expression-based copy-number inference at this sequencing depth.

Neither the supervised classifier nor CopyKAT responds to dosage
(Figure~\ref{fig:spike}b). The ikarus call rate on aneuploid cells (28.8\% at
12--18 events) is indistinguishable from its call rate on the diploid
companion cells in the same sample (30.9\%). Whatever ikarus is detecting, it
is not copy-number dosage; its calls track the transcriptional context of the
tissue, which is consistent with both its 100\% rate on Crohn's colonic
epithelium and its 0\% rate on other datasets. CopyKAT's response is erratic
and tissue-dependent: in spleen it calls 58.3\% of diploid companions and
94.9\% of cells carrying 12--18 events aneuploid, in liver 1.7\% and 0.0\%,
and in lung 14.9\% and 0.0\% --- that is, in two of three tissues the cells
with the largest imposed karyotypes are the ones it is least likely to call
aneuploid. Because CopyKAT emits a categorical call rather than a per-cell
score, it cannot be placed on the matched-false-positive-rate comparison
below.

The comparison that matters for practice is between methods held at the same
false-positive rate (Figure~\ref{fig:spike}c). For each rule, we measured its
false-call rate on the diploid companions and its detection rate on the
aneuploid cells, then thresholded the continuous score at exactly that
false-call rate and read off its detection rate. At every operating point,
ranking cells by the raw score detects substantially more aneuploid cells
than the rule that operates there: 23.9\% versus 5.1\% at 5.4\% false calls
(two-criterion rule), 65.0\% versus 21.5\% at 27.9\% false calls (forced
two-group split), and 48.8\% versus 28.8\% at 30.9\% false calls (ikarus),
all for cells carrying 12--18 arm events. The published rules therefore do
not extract information the score does not contain; they select a poor
operating point on it. The two-criterion rule is additionally uninformative
about burden --- its detection rate at 12--18 events (5.1\%) is no higher
than its detection rate averaged over all burdens (5.8\%).

\subsection{The false calls share a gene-density-driven pseudo-karyotype}

If false calls were random, they would be a power problem. They are not. We
compared the mean inferred copy-number profile of falsely called cells to
that of uncalled cells within each dataset, in 10-Mb genomic bins
(Figure~\ref{fig:mech}a). The resulting difference profile is highly
structured, with reproducible apparent losses and gains at specific loci ---
a pseudo-karyotype that a reader would naturally interpret as a shared clonal
event.

This profile is reproducible across tissues that share no biology: bin-level
difference profiles from independent datasets correlate with a median
Spearman $\rho$ of 0.54 (Figure~\ref{fig:mech}c). A genuine clonal event
cannot be shared by healthy kidney, healthy spleen and psoriatic skin.

The profile is predicted by local gene density (Figure~\ref{fig:mech}b).
Across 278 genomic bins, the apparent copy-number shift in falsely called
cells correlates negatively with the number of genes in the bin (Spearman
$\rho = -0.53$, $P = 2\times10^{-21}$): gene-dense regions appear deleted,
and gene-poor regions appear gained. This is the expected signature of a
compositional artefact. Sliding-window averaging of library-size-normalised
expression makes each window's value depend on relative, not absolute,
expression; cells whose expression distribution is broader or narrower than
the reference --- for reasons of cell state, RNA content, or capture
efficiency --- acquire a systematic, gene-density-dependent deviation across
the whole genome. Because that deviation is genome-wide and smooth, it
survives the smoothing steps copy-number callers apply, and because it is
shared by cells with similar technical properties, it clusters those cells
together and makes them look like a clone. We tested and rejected an
alternative explanation: clustered gene families (immunoglobulin, HLA,
histone, keratin and other locus-restricted families) do not explain the
profile as a family-level effect, although individual family loci appear
among the strongest bins.

\subsection{A calibrated decision procedure with a verifiable error rate}

The preceding results locate the problem in the decision step, not in the
copy-number score. We therefore replaced the decision step. The procedure
fits an empirical null distribution of the per-cell copy-number score on
tumour-free reference cells, conditional on cell type, residualising
sequencing depth --- because the score depends on depth, and cell types
differ in depth. A query cell is called aneuploid if its score exceeds the
$1-\alpha$ quantile of the null for its cell type. The procedure consumes any
per-cell score and does not modify the caller.

Under leave-one-tissue-out validation --- the null fitted on healthy tissues
excluding the held-out dataset --- the realized false-positive rate is at or
below the nominal target across the range $\alpha = 1\%$ to $10\%$
(Figure~\ref{fig:cal}a), in 13 of 14 datasets. At $\alpha = 5\%$ the median
realized rate is 0.03\% and the maximum 0.43\%, compared with 0.5--19.3\% for
the two-criterion rule on the same data (Figure~\ref{fig:cal}b). The
procedure is conservative rather than exact: a null pooled over ten tissues
has heavier tails than any single tissue, so the realized rate falls below
nominal. For an audit application this is the desired direction of error,
but users seeking an exact rate should fit the null closer to the query data.

Because cell-type labels are not always available, we also implemented a
label-free variant that assigns each cell to one of a small set of
transcriptional archetypes derived from reference centroids and conditions on
the archetype instead of the label. It behaves similarly at $\alpha \le 5\%$,
with more variability at $\alpha = 10\%$ (Figure~\ref{fig:cal}a).

The single failure is informative. On the psoriatic skin dataset, the
cross-tissue null gives a realized rate of 27.1\% --- worse than the
published rule it replaces. That dataset differs from the reference tissues
in protocol and study of origin, and the score distribution shifts
accordingly, so a null fitted elsewhere does not transfer. Refitting the null
on held-out donors from the same dataset removes the failure entirely: 27.1\%
$\rightarrow$ 0.03\%, with a maximum of 3.67\% across all datasets analysed
this way (Figure~\ref{fig:cal}c). The practical recommendation follows
directly: the null must be drawn from cells processed like the query cells,
and where a study contains multiple donors, held-out donors from the same
study are the appropriate source.

\section{Discussion}

Malignant-cell calling in single-cell RNA-seq is performed with tools whose
false-positive rate on tumour-free tissue has not been measured. We measured
it. Under published decision rules the rate is not small, not stable across
tissues, and not primarily a property of the biology: it is set by where an
essentially arbitrary threshold falls on a unimodal score.

The three findings compound. The score has limited per-cell discrimination
for realistic karyotype burdens; the rules applied to it operate at points
that waste even that discrimination; and the errors they make are
systematically structured rather than random, producing an apparent shared
karyotype driven by gene density. The last point is the most consequential
for interpretation, because a structured false-positive population does not
merely dilute a malignant compartment --- it adds a coherent, reproducible,
biologically meaningless signal to it, of exactly the kind that clonal
analyses are designed to detect.

We do not conclude that copy-number inference from single-cell expression is
without value. The continuous score carries real dosage information, and at
high karyotype burden it discriminates well. The failure is in the conversion
of that score into a binary label without a calibrated error rate. A
cell-type-conditional empirical null, costing seconds of computation and
requiring no change to the caller, converts an uncontrolled decision into one
whose error rate can be stated and verified. Its conservativeness --- and its
failure under protocol shift, which we report rather than conceal --- should
both be understood before it is applied.

Several practical recommendations follow. Report a false-positive rate,
measured on tumour-free cells processed identically, alongside any malignant
compartment. Prefer the continuous score to a binary label wherever the
downstream analysis permits. Where a binary label is required, set the
threshold from an explicit null rather than from a clustering outcome or a
fixed quantile. Treat an apparent shared karyotype among called cells as a
warning sign, and check it against gene density before interpreting it as
clonal. And do not transfer a null across protocols.

\section{Limitations}

This study was executed at reduced scale relative to its design, and the
reductions are material enough to state explicitly.

\textbf{Panel size.} Each dataset was subsampled to approximately 3{,}000
cells to make the full factorial of callers, reference modes and decision
rules tractable on a single workstation. The panel therefore covers 14
datasets and 36{,}528 cells, not the full atlases from which they are drawn.
Rates reported per dataset carry sampling uncertainty accordingly; pooled
rates are more stable than individual ones.

\textbf{Callers covered.} Two callers were run across the full panel:
infercnvpy (a Python implementation of the inferCNV algorithm) and ikarus.
The original R implementation of inferCNV~\cite{infercnv} was not run;
infercnvpy implements
the same sliding-window algorithm but is not bit-identical, and rates
reported here should not be attributed to the R package without verification.
CopyKAT was run across the full panel, but on a subsample of at most 1{,}200
cells per dataset (591--1{,}210 cells, drawn proportionally across cell
types) rather than the approximately 3{,}000 cells used for the other two
callers, because of its single-threaded runtime; its rates therefore carry
wider sampling uncertainty. CopyKAT was run in its default immune-reference
configuration only, not across the five decision rules, since it emits a
categorical call rather than a per-cell score. Several further callers in
common use --- CONICSmat~\cite{conics}, SCEVAN~\cite{scevan},
Numbat~\cite{numbat} and other methods that use allelic
information~\cite{casper,honeybadger} --- were not run at all. Numbat and related allele-aware methods are of particular interest
because the artefact we describe is compositional and would not be expected
to affect allelic imbalance in the same way; this is the most important
untested extension of the audit.

\textbf{Positive controls are in silico.} The aneuploid cells are real cells
with simulated arm-level dosage changes, not cells from tumours with
karyotypes established by an orthogonal assay. The simulation imposes a clean
multiplicative dosage effect and therefore, if anything, overstates
detectability relative to real tumour cells, where dosage compensation,
subclonal structure and focal events complicate the signal. The AUC values we
report should be read as upper bounds.

\textbf{No matched orthogonal ground truth.} We did not analyse a dataset
with matched single-cell DNA or bulk copy-number profiling, which would allow
false-positive and false-negative rates to be measured on genuine tumour
tissue rather than inferred from tumour-free tissue and simulation.

\textbf{Composition of the inflammatory-cohort datasets.} Of the four
non-healthy datasets, three contain cells annotated as diseased and one
(colon, ulcerative colitis cohort) contains only cells annotated as
non-diseased from that cohort. Conclusions about inflammation as a driver of
false calls rest on the former three.

\textbf{Calibration validation.} Leave-one-tissue-out validation uses
tumour-free tissue only; the realized rate reported is a false-positive rate,
not a false-discovery rate in the presence of true positives. The
donor-held-out variant was restricted to datasets with at least two donors
contributing at least 200 cells each.

\section{Methods}

\subsection{Data}

Healthy tissues were obtained from a multi-organ human cell reference
atlas~\cite{tabula} via the CELLxGENE Census, selecting primary 10x 3$'$ v3
data from adult donors. Inflammatory-cohort datasets were obtained from the
same source~\cite{kong2023,gao2021pso,smillie2019}. Each
dataset was subsampled to at most approximately 3{,}000 cells, with cell-type
proportions capped to avoid domination by a single type and all donors
retained. Genes were restricted to those with an assigned chromosome and
position in the annotation used by the callers. Per-dataset composition is
given in Supplementary Table S1.

\subsection{Copy-number inference}

infercnvpy 0.6.1~\cite{infercnvpy} was run per dataset on log-normalised expression with
100-gene sliding windows and a 10-gene step, in two modes: reference-free (no
reference category) and immune-reference, in which cells whose annotated type
matched a curated immune vocabulary (T, B, NK, plasma, myeloid, monocyte,
macrophage, dendritic, mast, granulocyte and their subtypes) were designated
the reference category. Per-cell scores were computed from the resulting
window matrix as the mean squared deviation (used throughout unless stated),
the mean absolute deviation, and the correlation to the mean profile of the
top-scoring decile.

CopyKAT 1.2.5~\cite{copykat} was run in R 4.4.3 under the same immune-reference
designation, with default parameters, on a per-dataset subsample of at most
1{,}200 cells drawn proportionally across annotated cell types. Cells
returned as \texttt{not.defined} are reported separately and excluded from
rate denominators unless stated.

ikarus 0.0.3~\cite{ikarus} was run with the published pan-cancer gene
signature and logistic-regression classifier, using AUCell
scoring~\cite{aucell} followed by
cell-cell-network propagation as implemented in the package.

\subsection{Decision rules}

Five rules were applied to the infercnvpy output.
(R1) \emph{Forced two-group split}: Leiden clustering~\cite{leiden} of the
CNV profile,
clusters merged into two groups, the group with the higher mean score
declared malignant.
(R1b) \emph{Highest-scoring cluster}: Leiden clustering at default
resolution, the single highest-mean-score cluster declared malignant.
(R2) \emph{Per-cell quantile}: cells above the 90th percentile of the
per-cell score declared malignant.
(R3) \emph{Reference SD}: cells whose score exceeds the reference-cell mean
by more than a fixed multiple of the reference standard deviation; requires a
reference and is undefined in reference-free mode.
(R4) \emph{Two-criterion (Tirosh)}~\cite{tirosh2016,puram2017}: cells
exceeding both a score threshold
and a threshold on correlation to the mean profile of the top-scoring cells.

\subsection{In-silico aneuploid spike-ins}

For each of three healthy tissues, cells were sampled at random and assigned
a karyotype consisting of $b$ chromosome arms drawn without replacement, each
designated gain or loss with equal probability. Counts for genes on an
affected arm were scaled by $1.5$ (gain) or $0.5$ (loss) before
re-normalisation, so that the resulting cell has the library size and
technical characteristics of a real cell with an altered relative expression
profile. Burdens $b \in \{1,2,3,5,8,12,18\}$ were generated. Untouched cells
from the same tissue serve as diploid companions. Truth labels and burdens
were held out of all scoring.

\subsection{Calibration}

Let $s_i$ be the per-cell copy-number score and $d_i$ the log total counts.
Within each cell type $t$ of the reference set, a regression of $s$ on $d$ is
fitted and the residual $r_i = s_i - \hat{s}(d_i)$ retained; the empirical
distribution of $r$ within type $t$ is the null $F_t$. A query cell of type
$t$ receives $p_i = 1 - F_t(r_i)$ and is called aneuploid if $p_i \le \alpha$.
Types with fewer than 100 reference cells are pooled into a global null. The
label-free variant replaces $t$ with the nearest reference centroid among a
small set of transcriptional archetypes, computed on marker-gene-restricted
expression.

Leave-one-tissue-out validation fits the null on all healthy tissues except
the held-out dataset and applies it to the held-out dataset. The
donor-held-out variant fits the null on all donors of a dataset except one
and applies it to the held-out donor, cycling over donors contributing at
least 200 cells.

\subsection{Attribution analysis}

Per-dataset difference profiles were computed as the mean window value of
falsely called cells minus that of uncalled cells, aggregated into 10-Mb
bins, retaining bins containing at least five windows. Gene counts per bin
were computed from the same annotation used by the caller. Cross-tissue
reproducibility was quantified as the Spearman correlation between bin-level
difference profiles of different datasets, over bins present in at least six
datasets. Clustered gene families were tested by comparing the mean
difference in bins overlapping family loci with the genome-wide distribution.

\subsection{Software and availability}

Analyses used Python 3.13 with scanpy 1.11.5~\cite{scanpy}, anndata 0.12.19, infercnvpy
0.6.1, ikarus 0.0.3, scikit-learn 1.9.1, numpy 2.4.6 and scipy 1.17.1, and R
4.4.3 with CopyKAT 1.2.5. All analysis code, result tables and figure scripts
accompany this manuscript.

\subsection{Data availability}

All datasets are public. The healthy-organ datasets are from the
multi-organ reference atlas~\cite{tabula} as distributed by the CELLxGENE
Census. The inflammatory-cohort datasets are CELLxGENE dataset
\texttt{333491b4-ca9b-4369-8706-38f7bb5bcd46} (Crohn's colon epithelial) and
\texttt{1ddca82d-2948-4dbe-8bba-be9de7115bb7} (Crohn's colon
immune)~\cite{kong2023},
\texttt{75eae05f-4f07-4518-a434-ea0a0d7ac616} (psoriasis
skin)~\cite{gao2021pso} and
\texttt{774d77e7-210d-4237-8ca8-8d2bb841a2e4} (ulcerative colitis cohort
colon)~\cite{smillie2019}. The per-dataset panel definition, including
download URLs, accompanies the analysis code.

\section*{Figures}

\begin{figure}[htbp]
\centering
\includegraphics[width=\textwidth]{FIG1}
\caption{\textbf{Malignant-cell callers assign tumour labels to tumour-free
tissue.}
(a) Percentage of cells called malignant in each of 14 tumour-free datasets
by infercnvpy under the two-criterion (Tirosh) rule, by CopyKAT, and by the
supervised classifier ikarus. Every call shown is a false positive. The shaded region
marks datasets from inflammatory-disease cohorts. inferCNV values use the
immune-reference mode where an immune reference was available and the
reference-free mode otherwise.
(b) Pooled false-call rate across the panel for four decision rules applied
to the same copy-number profiles, in reference-free and immune-reference
modes. The reference-SD rule is undefined without a reference (n.d.).
(c) Distribution of the per-cell copy-number score in healthy lung
(immune-reference mode). The dashed line is the reference-SD cutoff realized
on this dataset, which lands inside the bulk of the distribution and calls
6.3\% of cells. The score is unimodal: there is no malignant mode to separate,
so any cut selects a tail of an otherwise continuous distribution.}
\label{fig:fp}
\end{figure}

\begin{figure}[htbp]
\centering
\includegraphics[width=\textwidth]{FIG2}
\caption{\textbf{In-silico positive controls bound what the callers can
detect.}
(a) Area under the ROC curve for separating aneuploid from diploid cells
using the continuous copy-number score, as a function of the number of
arm-level events imposed, in three tissues.
(b) Percentage of cells called malignant by ikarus (solid) and CopyKAT
(dashed) as a function of imposed karyotype burden; burden 0 denotes
untouched diploid cells from the same sample. Neither caller responds to
dosage, and in liver and lung CopyKAT calls the highest-burden cells least
often.
(c) Detection of aneuploid cells carrying 12--18 arm events, comparing each
published rule (grey) with a threshold on the continuous score placed at
exactly that rule's false-call rate (purple). Labels give each rule's
false-call rate on diploid companion cells. Values are means over three
spike-in datasets.}
\label{fig:spike}
\end{figure}

\begin{figure}[htbp]
\centering
\includegraphics[width=\textwidth]{FIG3}
\caption{\textbf{False calls share a reproducible, gene-density-driven
pseudo-karyotype.}
(a) Mean inferred copy-number difference between falsely called and uncalled
cells, in 10-Mb bins, averaged across the tumour-free panel. Labels mark the
three strongest bins.
(b) Apparent copy-number shift against the number of genes per 10-Mb bin,
with linear fit. Gene-dense regions appear deleted.
(c) Spearman correlation between bin-level difference profiles of independent
datasets. Reproducibility across tissues that share no biology indicates a
technical rather than clonal origin.}
\label{fig:mech}
\end{figure}

\begin{figure}[htbp]
\centering
\includegraphics[width=\textwidth]{FIG4}
\caption{\textbf{A cell-type-conditional empirical null gives a verifiable
false-positive rate.}
(a) Realized versus nominal false-positive rate under leave-one-tissue-out
validation, for cell-type-conditioned and label-free archetype-conditioned
nulls. Mean $\pm$ s.d. across 13 datasets; the protocol-shifted psoriatic
skin dataset is excluded here and shown in (b) and (c). The dashed line is
the nominal rate.
(b) False-call rate per dataset for the two-criterion rule (grey) and for the
calibrated procedure at $\alpha = 5\%$ (blue). As in Fig.~\ref{fig:fp}a, the
rule rates use the immune-reference mode where an immune reference was
available and the reference-free mode otherwise (Crohn's epithelium).
(c) Effect of the null's provenance: a null fitted on other tissues versus
one fitted on held-out donors of the same dataset, at $\alpha = 5\%$,
datasets ordered by the former.}
\label{fig:cal}
\end{figure}

\begin{thebibliography}{99}
\small
\bibitem{tirosh2016} Tirosh I, et al. Dissecting the multicellular ecosystem of metastatic melanoma by single-cell RNA-seq. \emph{Science} 2016;352:189-196. doi:\href{https://doi.org/10.1126/science.aad0501}{10.1126/science.aad0501}
\bibitem{tirosh2016b} Tirosh I, et al. Single-cell RNA-seq supports a developmental hierarchy in human oligodendroglioma. \emph{Nature} 2016;539:309-313. doi:\href{https://doi.org/10.1038/nature20123}{10.1038/nature20123}
\bibitem{patel2014} Patel AP, et al. Single-cell RNA-seq highlights intratumoral heterogeneity in primary glioblastoma. \emph{Science} 2014;344:1396-1401. doi:\href{https://doi.org/10.1126/science.1254257}{10.1126/science.1254257}
\bibitem{infercnv} inferCNV of the Trinity CTAT Project. Broad Institute of MIT and Harvard. \url{https://github.com/broadinstitute/infercnv}
\bibitem{copykat} Gao R, et al. Delineating copy number and clonal substructure in human tumors from single-cell transcriptomes. \emph{Nature Biotechnology} 2021;39:599-608. doi:\href{https://doi.org/10.1038/s41587-020-00795-2}{10.1038/s41587-020-00795-2}
\bibitem{conics} Müller S, et al. CONICS integrates scRNA-seq with DNA sequencing to map gene expression to tumor sub-clones. \emph{Bioinformatics} 2018;34:3217-3219. doi:\href{https://doi.org/10.1093/bioinformatics/bty316}{10.1093/bioinformatics/bty316}
\bibitem{casper} Serin Harmanci A. CaSpER identifies and visualizes CNV events by integrative analysis of single-cell or bulk RNA-sequencing data. \emph{Nature Communications} 2020;11:89. doi:\href{https://doi.org/10.1038/s41467-019-13779-x}{10.1038/s41467-019-13779-x}
\bibitem{honeybadger} Fan J, et al. Linking transcriptional and genetic tumor heterogeneity through allele analysis of single-cell RNA-seq data. \emph{Genome Research} 2018;28:1217-1227. doi:\href{https://doi.org/10.1101/gr.228080.117}{10.1101/gr.228080.117}
\bibitem{numbat} Gao T, et al. Haplotype-aware analysis of somatic copy number variations from single-cell transcriptomes. \emph{Nature Biotechnology} 2022;41:417-426. doi:\href{https://doi.org/10.1038/s41587-022-01468-y}{10.1038/s41587-022-01468-y}
\bibitem{scevan} De Falco A, et al. A variational algorithm to detect the clonal copy number substructure of tumors from scRNA-seq data. \emph{Nature Communications} 2023;14:1074. doi:\href{https://doi.org/10.1038/s41467-023-36790-9}{10.1038/s41467-023-36790-9}
\bibitem{ikarus} Dohmen J, et al. Identifying tumor cells at the single-cell level using machine learning. \emph{Genome Biology} 2022;23:123. doi:\href{https://doi.org/10.1186/s13059-022-02683-1}{10.1186/s13059-022-02683-1}
\bibitem{tabula} The Tabula Sapiens Consortium. The Tabula Sapiens: A multiple-organ, single-cell transcriptomic atlas of humans. \emph{Science} 2022;376:eabl4896. doi:\href{https://doi.org/10.1126/science.abl4896}{10.1126/science.abl4896}
\bibitem{kong2023} Kong L, et al. The landscape of immune dysregulation in Crohn’s disease revealed through single-cell transcriptomic profiling in the ileum and colon. \emph{Immunity} 2023;56:444-458.e5. doi:\href{https://doi.org/10.1016/j.immuni.2023.01.002}{10.1016/j.immuni.2023.01.002}
\bibitem{gao2021pso} Gao Y, et al. Single cell transcriptional zonation of human psoriasis skin identifies an alternative immunoregulatory axis conducted by skin resident cells. \emph{Cell Death \& Disease} 2021;12:450. doi:\href{https://doi.org/10.1038/s41419-021-03724-6}{10.1038/s41419-021-03724-6}
\bibitem{smillie2019} Smillie CS, et al. Intra- and Inter-cellular Rewiring of the Human Colon during Ulcerative Colitis. \emph{Cell} 2019;178:714-730.e22. doi:\href{https://doi.org/10.1016/j.cell.2019.06.029}{10.1016/j.cell.2019.06.029}
\bibitem{infercnvpy} infercnvpy: infer copy number variation (CNV) from single-cell transcriptomics data. \url{https://github.com/icbi-lab/infercnvpy}
\bibitem{aucell} Aibar S, et al. SCENIC: single-cell regulatory network inference and clustering. \emph{Nature Methods} 2017;14:1083-1086. doi:\href{https://doi.org/10.1038/nmeth.4463}{10.1038/nmeth.4463}
\bibitem{leiden} Traag VA, Waltman L, van Eck NJ. From Louvain to Leiden: guaranteeing well-connected communities. \emph{Scientific Reports} 2019;9:5233. doi:\href{https://doi.org/10.1038/s41598-019-41695-z}{10.1038/s41598-019-41695-z}
\bibitem{puram2017} Puram SV, et al. Single-Cell Transcriptomic Analysis of Primary and Metastatic Tumor Ecosystems in Head and Neck Cancer. \emph{Cell} 2017;171:1611-1624.e24. doi:\href{https://doi.org/10.1016/j.cell.2017.10.044}{10.1016/j.cell.2017.10.044}
\bibitem{scanpy} Wolf FA, Angerer P, Theis FJ. SCANPY: large-scale single-cell gene expression data analysis. \emph{Genome Biology} 2018;19:15. doi:\href{https://doi.org/10.1186/s13059-017-1382-0}{10.1186/s13059-017-1382-0}
\end{thebibliography}

\end{document}
